\documentclass{article}
\usepackage[margin=1in]{geometry}
\usepackage{longtable}
\begin{document}

\textbf{Legacy note:} This split-by-person draft complements the canonical roadmap in \texttt{Documentation/SRS.tex}.\\

\section*{Phase 5 Plan - Person 2 (Person B)}

\textbf{Source of truth:} \texttt{Documentation/SRS.tex}, Section 11.5, ``Phase 5: Build Replay, Backtesting, and Validation''.\\
\textbf{Interpretation note:} Person 2 owns the empirical evaluation layer in Phase 5. The job is to decide whether replayed candidates and alerts are historically useful under point-in-time-safe validation, not just whether the replay machinery runs.\\
\textbf{Implementation note:} Person 2 should consume the replay outputs produced by Person 1 rather than reconstructing raw windows independently. Evaluation, paper-trading, and case studies must remain conservative and historically honest.

\textbf{Phase duration:} 3-4 weeks\\
\textbf{Phase goal:} Evaluate replayed historical outputs with conservative validation and produce defendable case studies.

\section*{Owned Deliverables (From SRS)}
\begin{enumerate}
\item Implement event-family holdout, category holdout, and time-split validation.
\item Build the conservative paper-trading simulator with bid/ask execution, slippage, size caps, fees, and near-resolution restrictions.
\item Define success metrics and failure metrics for candidate quality, alert usefulness, paper-trade edge, and lead time.
\item Produce at least three historical case studies that replay from raw data to alert decision.
\end{enumerate}

\section*{Locked Implementation Decisions}
\begin{itemize}
\item All evaluation must be point-in-time correct and may not use information that would not have been available at alert time.
\item Validation must include event-family holdout, category holdout, and time-split analysis rather than one pooled score.
\item Paper-trading assumptions must be conservative and explicitly documented, especially around spread crossing, fees, slippage, size limits, and near-resolution exclusions.
\item Success metrics must include failure-oriented views, not just favorable precision or return summaries.
\item Case studies must be reproducible from replay artifacts and explain both why an alert fired and how it would have been judged historically.
\end{itemize}

\section*{Concrete Execution Breakdown}
\begin{longtable}{p{0.28\linewidth}p{0.32\linewidth}p{0.30\linewidth}}
\textbf{Work Item} & \textbf{Artifact} & \textbf{Done Condition} \\
Validation framework & Event-family, category, and time-split validation outputs & Historical evaluation is segmented in a way that avoids optimistic leakage \\
Conservative simulator & Paper-trading simulator with realistic frictions & Replay outputs can be translated into conservative paper-trade outcomes \\
Metric package & Success and failure metrics for alerts and trades & Evaluation includes candidate quality, alert usefulness, lead time, and paper-trade edge \\
Historical case studies & At least three replay-to-alert narratives & Each case study links replay inputs, alert output, and historical interpretation in a defendable way \\
\end{longtable}

\section*{Execution Sequence (No Day Timeline)}
\textbf{Block 1: Validation design}
\begin{itemize}
\item Implement event-family holdout, category holdout, and time-split validation on top of replayed artifacts.
\item Define exactly which outputs are judged at the candidate, alert, and paper-trade levels.
\item Reject any evaluation path that leaks future public evidence or future market state into the judgment.
\end{itemize}

\textbf{Block 2: Conservative paper-trading simulator}
\begin{itemize}
\item Build a simulator that uses bid/ask execution rather than optimistic midpoint fills.
\item Add slippage, fees, size caps, and near-resolution restrictions.
\item Keep the simulator explainable enough that an academic reviewer can understand the assumptions without reverse-engineering the code.
\end{itemize}

\textbf{Block 3: Metrics and reporting}
\begin{itemize}
\item Define success metrics and failure metrics for candidate quality, alert usefulness, lead time, and paper-trade edge.
\item Produce one backtesting report template that can be reused across historical windows.
\item Surface both strong and weak regimes so the results are honest rather than cherry-picked.
\end{itemize}

\textbf{Block 4: Historical case studies and signoff pack}
\begin{itemize}
\item Produce at least three historical case studies from replayed raw data through alert decision.
\item Summarize why the system looked useful, where it failed, and what assumptions mattered most.
\item Assemble the historical validation pack and Gate 5 precondition evidence for shared signoff.
\end{itemize}

\section*{Handoffs to Person 1}
\begin{itemize}
\item Validation findings that expose replay gaps, missing fields, or reproducibility issues.
\item Simulator requirements that need stronger replay artifacts or additional stored metadata.
\item Metric definitions and report templates that should become canonical for later phases.
\item Historical windows that were especially informative, broken, or high-signal for future regression tests.
\item Case-study outputs that should be added to the shared Gate 5 evidence packet.
\end{itemize}

\section*{Gate 5 Checklist (Person 2 Ownership)}
\begin{itemize}
\item Validation includes event-family, category, and time-split holdouts.
\item Paper-trading assumptions are conservative, explicit, and reproducible.
\item Success and failure metrics are documented for candidates, alerts, lead time, and trades.
\item At least three historical case studies are complete and rerunnable from replay artifacts.
\item The historical validation pack is strong enough for shared Gate 5 precondition signoff.
\end{itemize}

\end{document}
